\chapter{Struktur Aljabar}\label{ch:b1:structures}

Aturan perhitungan yang sama terus muncul kembali: pada bilangan bulat, bilangan real,
bilangan kompleks, kelas \omterm{def:b1:arith:congruence}{kekongruenan}, dan sebentar lagi pada polinomial
(\cref{ch:b1:poly}), vektor dan matriks
(\cref{ch:b1:vspaces,ch:b1:matrices}). Aljabar menyarikan pola
yang berulang itu lalu menamainya: \emph{\omterm{def:b1:structures:group}{grup}}, \emph{\omterm{def:b1:structures:ring}{ring}}, \emph{\omterm{def:b1:structures:field}{lapangan}}.
Membuktikan sebuah fakta sekali, pada taraf strukturnya, berarti membuktikannya untuk
setiap contohnya sekaligus.

\section{Operasi biner}

\begin{definition}\label{def:b1:structures:law}
Sebuah \emph{operasi biner}\index{operasi biner} pada \omterm{def:b1:logic:sets}{himpunan} $E$ adalah
\omterm{def:b1:logic:map}{pemetaan} $E \times E \to E$, ditulis $(x, y) \mapsto x * y$. Ia bersifat
\emph{asosiatif} bila $(x*y)*z = x*(y*z)$ selalu berlaku, dan \emph{komutatif}
bila $x * y = y * x$ selalu berlaku. Unsur $e$ disebut \emph{identitas} bila
$e * x = x * e = x$ untuk setiap $x$; lalu $x'$ disebut \emph{invers} $x$
bila $x * x' = x' * x = e$.
\end{definition}

\begin{proposition}[Ketunggalan]\label{prop:b1:structures:unique}
Sebuah operasi mempunyai paling banyak satu identitas; dan pada operasi asosiatif
beridentitas, setiap unsur mempunyai paling banyak satu invers.
\end{proposition}

\begin{proof}
Jika $e$ dan $e'$ dua identitas: $e = e * e' = e'$. Jika $x'$ dan $x''$
membalikkan $x$: $x' = x' * e = x' * (x * x'') = (x' * x) * x'' = e * x''
= x''$.
\end{proof}

\section{Grup}

\begin{definition}[Grup]\label{def:b1:structures:group}
Sebuah \emph{grup}\index{grup} $(G, *)$ adalah \omterm{def:b1:logic:sets}{himpunan} yang dilengkapi operasi asosiatif
yang mempunyai identitas dan yang di dalamnya setiap unsur mempunyai invers. Grupnya
disebut \emph{abelian}\index{grup abelian} bila operasinya
komutatif.
\end{definition}

\begin{example}\label{ex:b1:structures:groups}
$(\Z, +)$, $(\Q, +)$, $(\R, +)$, $(\C, +)$; $(\Q^*, \times)$,
$(\R^*, \times)$, $(\C^*, \times)$, $(\mathbb{U}_n, \times)$ (\omterm{def:b1:complex:unity}{akar
satuan}, \cref{def:b1:complex:unity}); \omterm{def:b1:logic:sets}{himpunan} $\mathfrak{S}(E)$ berisi
bijeksi \omterm{def:b1:logic:sets}{himpunan} $E$ pada dirinya sendiri, terhadap komposisi --- yaitu
\emph{grup simetri}\index{grup simetri} $E$, yang tak \omterm{def:b1:structures:group}{abelian} begitu
$\abs E \geq 3$. Bukan \omterm{def:b1:structures:group}{grup}: $(\N, +)$ (tanpa invers),
$(\Z, \times)$ (hanya $\pm 1$ yang terbalikkan).
\end{example}

\begin{proposition}[Aturan perhitungan]\label{prop:b1:structures:rules}
Dalam \omterm{def:b1:structures:group}{grup} $G$ (ditulis secara perkalian, dengan identitas $e$):
\begin{enumerate}
  \item pencoretan: $ax = ay \implies x = y$ dan $xa = ya \implies
        x = y$;
  \item $(ab)^{-1} = b^{-1} a^{-1}$ dan $(a^{-1})^{-1} = a$;
  \item untuk $a, b \in G$, masing-masing persamaan $ax = b$ dan $xa = b$ mempunyai
        tepat satu penyelesaian ($x = a^{-1}b$, dan $x = b a^{-1}$).
\end{enumerate}
\end{proposition}

\begin{proof}
(1) Kalikan dengan $a^{-1}$ pada sisi yang sesuai, sambil memakai
keasosiatifannya. (2) $(b^{-1}a^{-1})(ab) = b^{-1}(a^{-1}a)b = b^{-1}b =
e$ dan simetris dengan itu; ketunggalan inversnya menutup buktinya; adapun butir
keduanya adalah \cref{prop:b1:structures:unique} yang diterapkan pada $a^{-1}$. (3)
Substitusikan lalu pakai (1) untuk ketunggalannya.
\end{proof}

\begin{example}[Simetri sebuah persegi panjang]\label{ex:b1:structures:klein}
Persegi panjang (yang bukan persegi) mempunyai tepat empat isometri pada
dirinya sendiri: identitas $e$, pencerminan pada sumbu mendatar $h$,
pencerminan pada sumbu tegak $v$, dan setengah putaran $r$ terhadap
pusatnya. Komposisi menjadikan \omterm{def:b1:logic:sets}{himpunan} beranggota empat ini sebuah \omterm{def:b1:structures:group}{grup}: setiap
unsurnya menjadi inversnya sendiri ($h^2 = v^2 = r^2 = e$), dan
hasil kali dua unsur tak identitas yang berbeda adalah unsur ketiganya
($hv = vh = r$: mencerminkan pada kedua sumbu sama dengan setengah putaran). Tabel
lengkapnya simetris, jadi \omterm{def:b1:structures:group}{grupnya} \omterm{def:b1:structures:group}{abelian} --- namun ia
\emph{bukan} \omterm{def:b1:structures:group}{grup} yang sama dengan \omterm{def:b1:structures:group}{grup} rotasi $\mathbb U_4$ pada
\cref{ex:b1:structures:order}: di sana $\iu$ berorde $4$, sedangkan
di sini setiap unsurnya berorde $\leq 2$. Jadi dua \omterm{def:b1:structures:group}{grup} berukuran sama
dapat mempunyai struktur perkalian yang sungguh-sungguh berbeda
--- gambar di bawah menampilkan kedua tabelnya berdampingan. \omterm{def:b1:structures:group}{Grup}
beranggota empat ini kembali sebagai $\{\pm1\} \times \{\pm1\}$, dan
\cref{exo:b1:structures:7} menjelaskan mengapa setiap \omterm{def:b1:structures:group}{grup} yang semua
kuadratnya sepele pastilah, seperti yang ini, \omterm{def:b1:structures:group}{abelian}.
\end{example}

\begin{omfigure}
\begin{tikzpicture}[scale=0.62]
  \foreach \i/\l in {0/$e$, 1/$\iu$, 2/$-1$, 3/$-\iu$} {
    \node[font=\small] at ({\i + 0.5}, 0.5) {\l};
    \node[font=\small] at (-0.5, {-\i - 0.5}) {\l};
  }
  \fill[omDef!25] (0,-1) rectangle ++(1,1);
  \fill[omDef!25] (3,-2) rectangle ++(1,1);
  \fill[omDef!25] (2,-3) rectangle ++(1,1);
  \fill[omDef!25] (1,-4) rectangle ++(1,1);
  \draw[gray] (0,-4) grid (4,0);
  \foreach \i/\j/\l in {0/0/$e$, 0/1/$\iu$, 0/2/$-1$, 0/3/$-\iu$,
    1/0/$\iu$, 1/1/$-1$, 1/2/$-\iu$, 1/3/$e$,
    2/0/$-1$, 2/1/$-\iu$, 2/2/$e$, 2/3/$\iu$,
    3/0/$-\iu$, 3/1/$e$, 3/2/$\iu$, 3/3/$-1$}
    \node[font=\scriptsize] at ({\j+0.5}, {-\i-0.5}) {\l};
  \begin{scope}[xshift=7.5cm]
    \foreach \i/\l in {0/$e$, 1/$h$, 2/$v$, 3/$r$} {
      \node[font=\small] at ({\i + 0.5}, 0.5) {\l};
      \node[font=\small] at (-0.5, {-\i - 0.5}) {\l};
    }
    \fill[omThm!25] (0,-1) rectangle ++(1,1);
    \fill[omThm!25] (1,-2) rectangle ++(1,1);
    \fill[omThm!25] (2,-3) rectangle ++(1,1);
    \fill[omThm!25] (3,-4) rectangle ++(1,1);
    \draw[gray] (0,-4) grid (4,0);
    \foreach \i/\j/\l in {0/0/$e$, 0/1/$h$, 0/2/$v$, 0/3/$r$,
      1/0/$h$, 1/1/$e$, 1/2/$r$, 1/3/$v$,
      2/0/$v$, 2/1/$r$, 2/2/$e$, 2/3/$h$,
      3/0/$r$, 3/1/$v$, 3/2/$h$, 3/3/$e$}
      \node[font=\scriptsize] at ({\j+0.5}, {-\i-0.5}) {\l};
  \end{scope}
\end{tikzpicture}

{\small Dua \omterm{def:b1:structures:group}{grup} beranggota empat: $\mathbb U_4 = \{e, \iu,
-1, -\iu\}$ (kiri) dan \omterm{def:b1:structures:group}{grup} persegi panjang (kanan), dengan
kedudukan identitasnya diarsir. Di kiri identitasnya berkelok
(satu unsur berorde $4$ membangkitkan segalanya); di kanan ia
memenuhi diagonalnya (setiap unsur berkuadrat $e$). Tak ada penamaan ulang
yang dapat mengubah tabel yang satu menjadi yang lain: jadi kedua \omterm{def:b1:structures:group}{grupnya} tidak
isomorfik.}
\end{omfigure}

\begin{definition}[Subgrup]\label{def:b1:structures:subgroup}
\omterm{def:b1:logic:sets}{Himpunan} bagian $H$ sebuah \omterm{def:b1:structures:group}{grup} $G$ disebut \emph{subgrup}\index{subgrup}
(ditulis $H \leq G$) bila ia memuat $e$, tertutup terhadap operasinya
dan terhadap pembalikannya. Maka $H$ sendiri sebuah \omterm{def:b1:structures:group}{grup}.

\emph{Kriteria:} $H \subseteq G$ yang tak kosong merupakan subgrup jika dan hanya
jika
\[
\forall x, y \in H, \quad x y^{-1} \in H .
\]
\end{definition}

\begin{proof}[Bukti kriterianya]
\omterm{def:b1:structures:subgroup}{Subgrup} jelas memenuhinya. Sebaliknya, misalkan $H \neq \emptyset$
memenuhinya, lalu pilih $x_0 \in H$. Maka $e = x_0 x_0^{-1} \in H$; untuk
$y \in H$, $y^{-1} = e\,y^{-1} \in H$; dan untuk $x, y \in H$, $xy =
x (y^{-1})^{-1} \in H$.
\end{proof}

\begin{example}\label{ex:b1:structures:subgroups}
$\mathbb{U}_n \leq (\C^*, \times)$: ia tak kosong, dan untuk $z, w \in
\mathbb{U}_n$ berlaku $(zw^{-1})^n = z^n (w^n)^{-1} = 1$. \omterm{def:b1:structures:subgroup}{Subgrup}
$(\Z, +)$ tepat berupa $n\Z$ (dibuktikan pada
\cref{thm:b1:arith:gcd}). Irisan \omterm{def:b1:structures:subgroup}{subgrup} selalu menjadi
\omterm{def:b1:structures:subgroup}{subgrup}, tetapi gabungannya hampir tak pernah demikian
(\cref{exo:b1:structures:6}).
\end{example}

\begin{remark}[Jebakan yang lazim dengan struktur]\label{rem:b1:structures:pitfalls}
\begin{enumerate}
  \item \emph{Ketertutupan terhadap operasinya saja belum cukup.} \omterm{def:b1:logic:sets}{Himpunan} $\N$
        tertutup terhadap penjumlahan di dalam $\Z$ dan memuat $0$, namun ia
        bukan \omterm{def:b1:structures:subgroup}{subgrup}: karena inversnya tak ada. Kriteria
        $xy^{-1} \in H$ menguji semuanya sekaligus --- tetapi hanya
        setelah $H \neq \emptyset$ diperiksa.
  \item \emph{Refleks tak \omterm{def:b1:structures:group}{abelian}.} Pada \omterm{def:b1:structures:group}{grup} yang umum, $(ab)^2 =
        abab$, yang sama dengan $a^2b^2$ hanya bila $a$ dan $b$ berkomutasi;
        demikian pula $(ab)^{-1} = b^{-1}a^{-1}$, dengan urutan terbalik.
        Setiap kesamaan yang diimpor dari aljabar sekolah harus
        diturunkan ulang dari aksiomanya atau ditandai sebagai kasus komutatif.
  \item \emph{\omterm{def:b1:structures:morphism}{Kernel} berbanding peta.} Di sini $\ker f$ hidup di
        \emph{sumbernya}, sedangkan $\operatorname{im} f$ di sasarannya;
        ``$f$ \omterm{def:b1:logic:inj}{injektif} jika dan hanya jika $\ker f$ sepele''
        (\cref{prop:b1:structures:kernel}) tak mempunyai analog dengan
        petanya ($\operatorname{im} f = G'$ adalah kesurjektifan).
  \item \emph{\omterm{def:b1:structures:ring}{Ring} bukan \omterm{def:b1:structures:group}{grup} terhadap $\times$.} Dalam sebuah \omterm{def:b1:structures:ring}{ring}, kebanyakan
        unsurnya tak harus terbalikkan, dan pencoretan oleh $a$
        menuntut $a$ menjadi unit atau \omterm{def:b1:structures:ring}{ringnya} menjadi \omterm{def:b1:structures:field}{daerah
        integral}: di $\Z/12\Z$, $\overline3\,\overline2 =
        \overline3\,\overline6$ padahal $\overline2 \neq \overline6$
        (\cref{ex:b1:structures:zncomputation}).
\end{enumerate}
\end{remark}

\begin{definition}[Morfisma grup]\label{def:b1:structures:morphism}
Misalkan $(G, *)$ dan $(G', \star)$ dua \omterm{def:b1:structures:group}{grup}. \omterm{def:b1:logic:map}{Pemetaan} $f \colon G \to G'$
disebut \emph{morfisma}\index{morfisma!grup} bila
\[
\forall x, y \in G, \qquad f(x * y) = f(x) \star f(y).
\]
Maka $f(e_G) = e_{G'}$ dan $f(x^{-1}) = f(x)^{-1}$. Adapun
\emph{kernel}\index{kernel} dan \emph{peta} $f$ adalah
\[
\ker f = f^{-1}(\{e_{G'}\}) \leq G,
\qquad
\operatorname{im} f = f(G) \leq G' .
\]
Morfisma yang \omterm{def:b1:logic:inj}{bijektif} disebut \emph{isomorfisma}; dan \omterm{def:b1:logic:map}{pemetaan} inversnya lalu
otomatis menjadi morfisma pula.
\end{definition}

\begin{proof}[Bukti pernyataannya]
Dari $f(e) = f(e * e) = f(e)\star f(e)$, mencoret $f(e)$ memberikan
$e_{G'} = f(e)$. Lalu $f(x)\star f(x^{-1}) = f(x x^{-1}) = e_{G'}$
mengenali $f(x^{-1})$ sebagai inversnya. \omterm{def:b1:structures:morphism}{Kernelnya}: $e \in \ker f$; jika
$x, y \in \ker f$, maka $f(xy^{-1}) = f(x)f(y)^{-1} = e$; jadi kriterianya
berlaku. Petanya: kriteria yang sama dengan $f(x)f(y)^{-1} = f(xy^{-1})$.
Invers sebuah isomorfisma: untuk $u, v \in G'$, tulis $u = f(x)$, $v =
f(y)$; maka $f^{-1}(u \star v) = f^{-1}(f(xy)) = xy =
f^{-1}(u) f^{-1}(v)$.
\end{proof}

\begin{proposition}[Keinjektifan lewat kernelnya]\label{prop:b1:structures:kernel}
Sebuah \omterm{def:b1:structures:morphism}{morfisma} \omterm{def:b1:structures:group}{grup} $f$ bersifat \omterm{def:b1:logic:inj}{injektif} jika dan hanya jika $\ker f =
\{e\}$.
\end{proposition}

\begin{proof}
Jika $f$ \omterm{def:b1:logic:inj}{injektif}, maka $\ker f$ hanya dapat memuat satu \omterm{def:b1:logic:map}{prapeta}
$e_{G'}$, yaitu $e$. Sebaliknya, jika $\ker f = \{e\}$ dan $f(x) =
f(y)$, maka $f(xy^{-1}) = f(x) f(y)^{-1} = e_{G'}$, sehingga $xy^{-1} = e$,
yakni $x = y$.
\end{proof}

\begin{example}\label{ex:b1:structures:expmorphism}
$\exp \colon (\R, +) \to (\R_+^*, \times)$ adalah \omterm{def:b1:structures:morphism}{morfisma}
(karena $\eu^{x+y} = \eu^x \eu^y$) yang \omterm{def:b1:logic:inj}{bijektif}
(\cref{prop:b1:functions:expln}): jadi struktur penjumlahan dan
perkaliannya isomorfik --- itulah alasan historis keberadaan
logaritma. \omterm{def:b1:structures:morphism}{Morfisma} lain: $\theta \mapsto \eu^{\iu\theta}$ dari
$(\R, +)$ pada lingkaran satuan $(\mathbb{U}, \times)$, dengan \omterm{def:b1:structures:morphism}{kernel}
$2\pi\Z$.
\end{example}

\begin{example}[Morfisma tanda]\label{ex:b1:structures:signmorphism}
\omterm{def:b1:logic:map}{Pemetaan} $s \colon (\R^*, \times) \to (\{\pm1\}, \times)$ yang mengirim
$x$ ke tandanya adalah \omterm{def:b1:structures:morphism}{morfisma}: karena tanda sebuah hasil kali sama dengan
hasil kali tandanya. \omterm{def:b1:structures:morphism}{Kernelnya} adalah $\intoo0{+\infty}$ (sebuah
\omterm{def:b1:structures:subgroup}{subgrup}, sesuai janji \cref{def:b1:structures:morphism}), dan
petanya seluruh $\{\pm1\}$: jadi \omterm{def:b1:logic:inj}{surjektif}, tetapi sangat tak \omterm{def:b1:logic:inj}{injektif}.
Dua pelajaran umum dalam wujud mini. Pertama, sebuah \omterm{def:b1:structures:morphism}{morfisma} boleh saja meremukkan
informasi: $s$ tak mengingat apa pun dari $x$ selain satu bit, dan itulah
keutamaannya --- argumen tanda tepat merupakan perhitungan
yang melewati $s$. Kedua, \omterm{def:b1:structures:morphism}{morfisma} ke $\{\pm1\}$ adalah
``invarian'' yang paling sederhana: tanda \omterm{def:b1:counting:objects}{permutasi}, yang dibangun pada
soal akhir pekan bab ini, adalah gejala yang sama pada
\omterm{def:b1:structures:group}{grup} $\mathfrak S_n$, dan semua argumen paritas yang digerakkannya
turun lewat \omterm{def:b1:structures:morphism}{morfisma} bernilai dua semacam itu.
\end{example}

\begin{definition}[Pangkat, orde sebuah unsur]\label{def:b1:structures:order}
Dalam \omterm{def:b1:structures:group}{grup} $G$ (dengan notasi perkalian), tulis $x^0 = e$, $x^{k+1} =
x^k x$ dan $x^{-k} = (x^k)^{-1}$ untuk $k \in \N$; maka $x^{k+l} = x^k
x^l$ untuk setiap $k, l \in \Z$, sehingga $k \mapsto x^k$ adalah \omterm{def:b1:structures:morphism}{morfisma} $(\Z,
+) \to G$ yang petanya $\langle x \rangle = \{x^k : k \in \Z\}$ merupakan
\omterm{def:b1:structures:subgroup}{subgrup}, yaitu \omterm{def:b1:structures:subgroup}{subgrup} yang \emph{dibangkitkan} oleh $x$. Adapun
\emph{orde}\index{orde sebuah unsur} $x$ adalah $m \geq
1$ terkecil dengan $x^m = e$ bila ada (lalu $\langle x\rangle = \{e, x,
\dots, x^{m-1}\}$ beranggota tepat $m$ unsur, dan $x^k = e \iff m
\mid k$), dan $\infty$ bila tidak ada.
\end{definition}

\begin{example}\label{ex:b1:structures:order}
Dalam $(\C^*, \times)$: $\iu$ berorde $4$, dengan $\langle \iu \rangle
= \{1, \iu, -1, -\iu\} = \mathbb{U}_4$; lebih umum $\omega =
\eu^{2\iu\pi/n}$ berorde $n$ dan $\langle\omega\rangle =
\mathbb{U}_n$. Dalam $(\Z, +)$, setiap $x \neq 0$ berorde tak hingga.
Mengapa klaim pada definisinya berlaku: jika $x$ berorde $m$, bagilah
sebarang $k$ dengan $m$ ($k = mq + r$, $0 \leq r < m$,
\cref{thm:b1:arith:division}): maka $x^k = (x^m)^q x^r = x^r$, jadi
pangkatnya berdaur dengan periode $m$, unsur yang didaftar itu berbeda
dua-dua menurut keminimalan $m$, dan $x^k = e$ memaksa $r = 0$. \omterm{def:b1:structures:order}{Orde}
\omterm{def:b1:counting:objects}{permutasi} dihitung pada soal akhir pekan di bawah.
\end{example}

\begin{example}[Orde di dalam $\mathbb U_{12}$]\label{ex:b1:structures:orderU12}
Berapa \omterm{def:b1:structures:order}{orde} $\omega^k$ di $\mathbb U_n$, untuk $\omega =
\eu^{2\iu\pi/n}$? Kita punya $(\omega^k)^m = 1$ jika dan hanya jika $n \mid km$, dan
dengan menulis $d = \gcd(n, k)$, $n = dn'$, $k = dk'$ dengan $\gcd(n', k')
= 1$: $n \mid km \iff n' \mid k'm \iff n' \mid m$ (menurut lema
Gauss, \cref{thm:b1:arith:gauss}). Nilai $m \geq 1$ terkecil semacam itu adalah
$n' = \frac{n}{\gcd(n,k)}$. Di $\mathbb U_{12}$ misalnya,
$\omega^8$ berorde $\frac{12}{\gcd(12,8)} = 3$ (memang
$\omega^8 = \eu^{4\iu\pi/3} \in \mathbb U_3$), sedangkan $\omega^5$
berorde $12$: ia membangkitkan seluruh \omterm{def:b1:structures:group}{grupnya}, meskipun ia bukan
pembangkit yang ``baku''. Mencacah pembangkitnya --- yaitu $k$
dengan $\gcd(k, n) = 1$ --- memulihkan cacahan \omterm{cor:b1:arith:bezout}{saling prima} pada
\cref{ex:b1:counting:coprime}: jadi teori \omterm{def:b1:structures:group}{grup} dan pencacahan bertemu.
\end{example}

\section{Ring dan lapangan}

\begin{definition}[Ring]\label{def:b1:structures:ring}
Sebuah \emph{ring}\index{ring} $(A, +, \times)$ adalah \omterm{def:b1:logic:sets}{himpunan} dengan dua operasi
sedemikian sehingga: $(A, +)$ \omterm{def:b1:structures:group}{grup abelian} (dengan identitas $0$); $\times$
asosiatif dengan identitas $1$; dan $\times$ distributif terhadap $+$ pada
kedua sisinya. Ringnya disebut \emph{komutatif} bila $\times$ komutatif. Sebuah
unsur $a$ disebut \emph{terbalikkan} (sebuah \emph{unit}) bila $ab = ba = 1$
untuk suatu $b$; dan semua unitnya membentuk \omterm{def:b1:structures:group}{grup} $(A^\times, \times)$.
\end{definition}

\begin{proof}[Bukti bahwa unitnya membentuk grup]
Ketertutupannya: jika $a, a'$ unit dengan invers $b, b'$, maka
\[
(aa')(b'b) = a(a'b')b = a\,1\,b = ab = 1,
\qquad (b'b)(aa') = 1
\]
secara simetris, jadi $aa'$ sebuah unit. Unsur $1$ adalah unit (yang
menjadi inversnya sendiri), keasosiatifannya diwarisi dari $A$, dan
invers $b$ sebuah unit $a$ sendiri merupakan unit (dengan invers $a$).
Jadi $(A^\times, \times)$ memenuhi semua aksioma \omterm{def:b1:structures:group}{grup}. Setiap \omterm{def:b1:structures:group}{grup}
dalam buku ini yang tidak dibangun dari \omterm{def:b1:counting:objects}{permutasi} muncul dengan cara ini:
$\Q^* = \Q^\times$, $\R^*$, $\C^*$, unit $\Z/n\Z$ di bawah,
dan nanti matriks yang terbalikkan (\cref{ch:b1:matrices}).
\end{proof}

\begin{example}\label{ex:b1:structures:rings}
$\Z, \Q, \R, \C$ adalah \omterm{def:b1:structures:ring}{ring} komutatif; $\Z^\times = \{1, -1\}$,
$\Q^\times = \Q^*$. Menyusul kemudian: \omterm{def:b1:structures:ring}{ring} polinomial $K[X]$
(\cref{ch:b1:poly}), \omterm{def:b1:structures:ring}{ring} matriks (yang tak komutatif,
\cref{ch:b1:matrices}), dan $\Z/n\Z$ di bawah. Dalam setiap \omterm{def:b1:structures:ring}{ring}, $0 \times
a = 0$ (dari sifat distributifnya: $0a = (0+0)a = 0a + 0a$), dan
$(-1)a = -a$.
\end{example}

\begin{example}[Idempoten: gejala baru dalam ring baru]\label{ex:b1:structures:idempotents}
Di $\Z$, persamaan $x^2 = x$, yakni $x(x - 1) = 0$, hanya mempunyai
penyelesaian $0$ dan $1$. Di $\Z/6\Z$, dengan menguji semua kelasnya:
$\overline0^2 = \overline0$, $\overline1^2 = \overline1$,
$\overline3^2 = \overline9 = \overline3$ dan $\overline4^2 =
\overline{16} = \overline4$ --- jadi ada \emph{empat} idempoten. Kedua
yang eksotis itu datang dari pembagi nol: $\overline3\,(\overline3 -
\overline1) = \overline3 \times \overline2 = \overline6 =
\overline0$ tanpa satu pun faktornya nol. Perhitungan semacam itu mengalibrasi
gerak hati kita: fakta yang sudah dikenal tentang persamaan bertahan di
\omterm{def:b1:structures:field}{daerah integral} dan \omterm{def:b1:structures:field}{lapangan}, tetapi \omterm{def:b1:structures:ring}{ring} yang umum dapat dan memang
berperilaku lain --- lihat pula \omterm{def:b1:structures:ring}{ring} Boole pada
\cref{exo:b1:structures:10}, yang di sana \emph{setiap} unsurnya
idempoten.
\end{example}

\begin{proposition}[Teorema binomial dalam ring komutatif]\label{prop:b1:structures:binomial}
Jika $a, b$ unsur sebuah \omterm{def:b1:structures:ring}{ring} komutatif (lebih umum lagi, jika
$ab = ba$), maka untuk $n \in \N$:
\[
(a+b)^n = \sum_{k=0}^n \binom nk a^k b^{n-k},
\qquad
a^n - b^n = (a - b) \sum_{k=0}^{n-1} a^k b^{\,n-1-k} .
\]
\end{proposition}

\begin{proof}
Bukti \cref{thm:b1:counting:binomial} dan bukti kesamaan geometrinya
hanya memakai keasosiatifan, kekomutatifan kedua unsurnya,
dan sifat distributifnya --- jadi keduanya berlaku kata demi kata.
\end{proof}

\begin{example}[Teorema binomial dalam ring yang tak dikenal]\label{ex:b1:structures:binomialring}
Dua hasil cepat dari keumumannya. Di $\Z/p\Z$ (dengan $p$ prima),
\omterm{def:b1:counting:objects}{koefisien binomial} di tengahnya lenyap
(yaitu langkah pertama \cref{thm:b1:arith:fermat}), sehingga teoremanya
runtuh menjadi \emph{mimpi mahasiswa baru}
\[
(a + b)^p = a^p + b^p \qquad \text{di } \Z/p\Z ,
\]
yang di sana merupakan kesamaan yang sungguhan, betapapun ia tampak jahat di atas $\R$.
Dan di dalam sebarang \omterm{def:b1:structures:ring}{ring} komutatif yang memuat unsur $\varepsilon$
dengan $\varepsilon^2 = 0$, teoremanya terpenggal:
$(a + \varepsilon)^n = a^n + n\,a^{n-1}\varepsilon$, karena semua suku
yang lebih tinggi memikul faktor $\varepsilon^2 = 0$. Koefisien
$n\,a^{n-1}$ pada $\varepsilon$ adalah turunan $x^n$ ---
dan itu bukan kebetulan, melainkan isyarat pertama bahwa turunan sama aljabarnya
dengan analisisnya (bandingkan turunan formal pada
\cref{ch:b1:poly}).
\end{example}

\begin{definition}[Daerah integral, lapangan]\label{def:b1:structures:field}
\omterm{def:b1:structures:ring}{Ring} komutatif $A \neq \{0\}$ disebut \emph{daerah
integral}\index{daerah integral} bila ia tidak mempunyai pembagi nol: yakni $ab = 0
\implies a = 0$ atau $b = 0$. Ia disebut \emph{lapangan}\index{lapangan} bila
setiap unsurnya yang tak nol terbalikkan. Setiap lapangan merupakan daerah
integral ($ab = 0$ dan $a \neq 0$ memberikan $b = a^{-1}ab = 0$).
\end{definition}

\begin{example}
$\Q$, $\R$, $\C$ adalah \omterm{def:b1:structures:field}{lapangan}; sedangkan $\Z$ \omterm{def:b1:structures:field}{daerah integral} tetapi bukan
\omterm{def:b1:structures:field}{lapangan}. Dalam \omterm{def:b1:structures:field}{daerah integral}, pencoretan berlaku untuk $\times$: $ab =
ac$ dan $a \neq 0$ mengakibatkan $b = c$.
\end{example}

\section{Ring $\Z/n\Z$}

\begin{definition}\label{def:b1:structures:zn}
Tetapkan $n \in \N^*$. Kelas \omterm{def:b1:arith:congruence}{kekongruenan} modulo $n$
(\cref{ex:b1:logic:congruence}) membentuk \omterm{def:b1:logic:sets}{himpunan}
$\Z/n\Z$\index{Z/nZ@$\Z/n\Z$} beranggota $n$ unsur, ditulis $\overline 0,
\overline 1, \dots, \overline{n-1}$. Operasinya
\[
\overline a + \overline b = \overline{a + b},
\qquad
\overline a \times \overline b = \overline{ab}
\]
terdefinisi dengan baik --- kelas hasilnya tidak bergantung pada
wakilnya, justru karena \omterm{def:b1:arith:congruence}{kekongruenannya} serasi dengan $+$
dan $\times$ (\cref{def:b1:arith:congruence}) --- dan menjadikan $\Z/n\Z$
sebuah \omterm{def:b1:structures:ring}{ring} komutatif.
\end{definition}

\begin{theorem}[Unit $\Z/n\Z$; lapangan $\Z/p\Z$]\label{thm:b1:structures:zp}
\begin{enumerate}
  \item $\overline a$ terbalikkan di $\Z/n\Z$ jika dan hanya jika
        $\gcd(a, n) = 1$.
  \item $\Z/n\Z$ merupakan \omterm{def:b1:structures:field}{lapangan} jika dan hanya jika $n$ prima.
\end{enumerate}
\end{theorem}

\begin{proof}
(1) adalah \cref{prop:b1:arith:invmod} yang ditulis ulang dengan kelas.

(2) Jika $n = p$ prima, setiap $\overline a \neq \overline 0$ memenuhi $p
\nmid a$, jadi $\gcd(a, p) = 1$: sehingga terbalikkan menurut (1) --- yaitu sebuah \omterm{def:b1:structures:field}{lapangan}. Jika $n =
ab$ dengan $1 < a, b < n$, maka $\overline a\, \overline b = \overline
n = \overline 0$ dengan $\overline a, \overline b \neq \overline 0$:
jadi ada pembagi nol, sehingga ia bahkan bukan \omterm{def:b1:structures:field}{daerah integral}; adapun $n = 1$ memberikan
\omterm{def:b1:structures:ring}{ring} nol, yang tersingkir.
\end{proof}

\begin{example}[Berapa banyak akar kuadrat dari $1$?]\label{ex:b1:structures:sqrtone}
Selesaikan $x^2 = \overline 1$ di $\Z/8\Z$ dan di $\Z/7\Z$. Dengan menguji
kedelapan kelas modulo $8$: $1^2 = 1$, $3^2 = 9 \equiv 1$, $5^2 = 25
\equiv 1$, $7^2 = 49 \equiv 1$ --- jadi ada \emph{empat} penyelesaian
$\{\overline1, \overline3, \overline5, \overline7\}$, padahal
polinomial $X^2 - 1$ berderajat $2$. Sebaliknya, di \omterm{def:b1:structures:field}{lapangan} $\Z/7\Z$,
$x^2 = \overline1$ berarti $(x - \overline1)(x +
\overline1) = \overline0$, dan \omterm{def:b1:structures:field}{lapangan} tak mempunyai pembagi nol: sehingga $x =
\pm\overline1$, dengan dua penyelesaian saja. Kegagalan modulo $8$ itu dapat
dilacak: $(3-1)(3+1) = 2 \times 4 = 8 \equiv 0$ tanpa satu pun
faktornya lenyap. Moralnya: kaidah yang sudah dikenal ``persamaan berderajat
$d$ mempunyai paling banyak $d$ akar'' adalah teorema tentang
\emph{\omterm{def:b1:structures:field}{daerah integral}} (\cref{cor:b1:poly:nroots} membuktikannya atas
\omterm{def:b1:structures:field}{lapangan}); dan di dalam \omterm{def:b1:structures:ring}{ring} yang mempunyai pembagi nol ia gugur diam-diam --- justru
itulah sebabnya bukti pemasangan pada teorema Wilson
(\cref{exo:b1:arith:11}) menuntut $p$ prima.
\end{example}

\begin{example}[Berhitung di $\Z/n\Z$]\label{ex:b1:structures:zncomputation}
Di $\Z/12\Z$: unitnya adalah $\overline 1, \overline 5, \overline 7,
\overline{11}$ (yaitu kelas yang \omterm{cor:b1:arith:bezout}{saling prima} dengan $12$), dan masing-masing menjadi
inversnya sendiri ($5^2 = 25 \equiv 1$, $7^2 = 49 \equiv 1$, $11^2 = 121 \equiv
1$). Persamaan $\overline 3\, x = \overline 6$ mempunyai \emph{tiga}
penyelesaian ($x \in \{\overline 2, \overline 6, \overline{10}\}$):
karena tanpa keterbalikan tak ada pencoretan. Sebaliknya di $\Z/11\Z$,
setiap persamaan $\overline a x = \overline b$ dengan $\overline a \neq
\overline 0$ mempunyai tepat satu penyelesaian.
\end{example}

\begin{example}[Aksioma grup sebagai izin menyelesaikan]\label{ex:b1:structures:solving}
Dalam \omterm{def:b1:structures:group}{grup} $\bigl((\Z/7\Z)^*, \times\bigr)$, selesaikan $\overline
3\,x = \overline 5$. Menurut \cref{prop:b1:structures:rules}~(3)
penyelesaiannya ada, tunggal, dan sama dengan $\overline3^{-1}\,
\overline5$; karena $\overline3 \times \overline5 = \overline{15}
= \overline1$, invers $\overline 3$ adalah $\overline 5$, sehingga
\[
x = \overline5 \times \overline5 = \overline{25} = \overline4,
\qquad\text{periksa: } \overline3 \times \overline4 =
\overline{12} = \overline5 .
\]
Yang penting di sini bukan jawabannya melainkan jaminannya: dalam sebuah \omterm{def:b1:structures:group}{grup},
setiap persamaan semacam itu terselesaikan secara tunggal \emph{sebelum} ada
perhitungan, sehingga prosedur penyelesaiannya tak akan pernah tersandung ``tak ada
penyelesaian'' atau ``ada beberapa''. Bandingkan $\overline3\,x = \overline6$
di $\Z/12\Z$ di atas, yang di sana jaminannya gugur --- mengetahui struktur
mana yang sedang kita huni berarti mengetahui apa yang boleh kita anggap sudah pasti.
\end{example}

\begin{example}[Hasil kali langsung]\label{ex:b1:structures:product}
Jika $G$ dan $H$ dua \omterm{def:b1:structures:group}{grup}, maka \omterm{def:b1:logic:sets}{himpunan} hasil kali $G \times H$ dengan
operasi komponen demi komponen $(g, h)(g', h') = (gg', hh')$ merupakan \omterm{def:b1:structures:group}{grup}:
aksiomanya diperiksa koordinat demi koordinat, dengan identitas
$(e_G, e_H)$ dan invers $(g^{-1}, h^{-1})$. \omterm{def:b1:structures:order}{Ordenya} berpadu lewat
KPK: $(g, h)^m = (g^m, h^m)$ menjadi identitas jika dan hanya jika \omterm{def:b1:structures:order}{orde}
$g$ maupun \omterm{def:b1:structures:order}{orde} $h$ \omterm{def:b1:arith:divides}{membagi} $m$. Jadi di $\Z/2\Z
\times \Z/2\Z$ (secara aditif) setiap unsur tak nol berorde $2$
--- dan ini tepat \omterm{def:b1:structures:group}{grup} persegi panjang pada
\cref{ex:b1:structures:klein} dalam koordinat --- sedangkan $\Z/4\Z$
mempunyai unsur berorde $4$: yaitu bukti kedua, yang bebas perhitungan,
bahwa kedua \omterm{def:b1:structures:group}{grup} berukuran $4$ itu tidak isomorfik (karena
isomorfisma mengawetkan \omterm{def:b1:structures:order}{orde}). Hasil kali adalah cara termudah
memproduksi \omterm{def:b1:structures:group}{grup} baru dari \omterm{def:b1:structures:group}{grup} lama, dan bidang $\R^2 = \R \times
\R$ pada \cref{ch:b1:vspaces} adalah wujud konstruksi itu yang paling
penting.
\end{example}

\begin{remark}[Fermat, secara struktural]
Dalam \omterm{def:b1:structures:field}{lapangan} $\Z/p\Z$, kelas yang tak nol membentuk \omterm{def:b1:structures:group}{grup} perkalian
beranggota $p - 1$ unsur, dan teorema kecil Fermat
(\cref{thm:b1:arith:fermat}) mengatakan: setiap unsur $x$ \omterm{def:b1:structures:group}{grup} ini
memenuhi $x^{p-1} = \overline 1$. Ini wujud sebuah fakta umum
tentang \omterm{def:b1:structures:group}{grup} hingga (teorema Lagrange), yang dibuktikan pada tahun
kedua; dan bukti pemasangan pada teorema Wilson
(\cref{exo:b1:arith:11}) sudah bercita rasa teori \omterm{def:b1:structures:group}{grup} itu.
\end{remark}

\begin{remark}[Selingan: apa yang dibeli abstraksi]\label{rem:b1:structures:interlude}
Wajar bertanya apa yang diperoleh dengan membuktikan, katakanlah,
\cref{prop:b1:structures:unique} untuk operasi yang abstrak alih-alih
untuk bilangan. Jawabannya adalah daya ungkit. Argumen dua baris itu kini
mencakup, sekaligus: invers fungsi terhadap komposisi
(\cref{thm:b1:logic:inverse}, yang bukti ketunggalannya diulanginya
kata demi kata), invers modulo $n$
(\cref{prop:b1:arith:invmod}), invers bilangan real tak nol, invers
unit dalam sebarang \omterm{def:b1:structures:ring}{ring}, dan --- tanpa perlu dilihat lebih dulu --- invers
matriks yang terbalikkan pada \cref{ch:b1:matrices}, yang di sana ketunggalan $A^{-1}$
tak akan menuntut satu baris bukti pun. Penghematan yang sama berlaku bagi
\cref{prop:b1:structures:kernel} (satu kriteria keinjektifan,
yang dipakai ulang untuk \omterm{def:b1:logic:map}{pemetaan} linear pada \cref{ch:b1:linmaps}) dan bagi
kriteria \omterm{def:b1:structures:subgroup}{subgrupnya}. Abstraksi di sini bukanlah keumuman demi keumuman
itu sendiri: ia adalah penolakan membuktikan lema yang sama lima kali
dengan lima nama. Ongkosnya --- yaitu melacak aksioma mana yang benar-benar
dipakai setiap \omterm{def:b1:logic:statement}{pernyataan} --- justru itulah yang dilatih latihan
pada bab ini.
\end{remark}

\begin{remark}[Di mana bab ini dipakai]\label{rem:b1:structures:whereused}
Kosakata bab ini adalah tata bahasa bagi selebihnya jilid
ini. \omterm{def:b1:structures:ring}{Ring} dan \omterm{def:b1:structures:field}{lapangan} menata \cref{ch:b1:poly} ($K[X]$ adalah
\omterm{def:b1:structures:ring}{ring} yang meniru $\Z$) dan \cref{ch:b1:fractions} ($K(X)$ adalah
\omterm{def:b1:structures:field}{lapangan} pecahannya); ruang vektor (\cref{ch:b1:vspaces}) adalah
\omterm{def:b1:structures:group}{grup abelian} yang di atasnya sebuah \omterm{def:b1:structures:field}{lapangan} bekerja; matriks
(\cref{ch:b1:matrices}) membentuk \omterm{def:b1:structures:ring}{ring} pertama dalam jilid ini yang
sungguh-sungguh tak komutatif, dan unsurnya yang terbalikkan membentuk \omterm{def:b1:structures:group}{grup} yang
telaahnya adalah aljabar linear itu sendiri. \omterm{def:b1:structures:morphism}{Morfisma} dan \omterm{def:b1:structures:morphism}{kernel} kembali sebagai
\omterm{def:b1:logic:map}{pemetaan} linear dan ruang nol pada \cref{ch:b1:linmaps} ---
\cref{prop:b1:structures:kernel} \emph{adalah} kriteria keinjektifan
pada bab itu, yang dibuktikan sekali untuk selamanya di sini. Adapun
\omterm{ex:b1:structures:groups}{grup simetri}, bintang soal akhir pekan di bawah, memasok
tanda yang di atasnya determinan dibangun pada \cref{ch:b1:det}.
\end{remark}

\section{Latihan}

\begin{exercise}[$\star$]\label{exo:b1:structures:1}
Pada $E = \R \setminus \{1\}$, definisikan $x * y = x + y - xy$. Buktikan bahwa
$(E, *)$ adalah \omterm{def:b1:structures:group}{grup abelian}. \emph{(Kenali identitasnya dan
invers $x$; periksa ketertutupannya: mengapa $x * y \neq 1$?)}
\end{exercise}

\begin{exercise}[$\star$]\label{exo:b1:structures:2}
Manakah di antara berikut ini yang merupakan \omterm{def:b1:structures:group}{grup}?
\begin{enumerate}
  \item $(\intoo{0}{+\infty}, \times)$;
  \item $(\{-1, 0, 1\}, +)$;
  \item $(\Q^*, \times)$;
  \item \omterm{def:b1:logic:sets}{himpunan} bilangan bulat ganjil terhadap penjumlahan.
\end{enumerate}
\end{exercise}

\begin{exercise}[$\star$]\label{exo:b1:structures:3}
Tuliskan tabel komposisi \omterm{ex:b1:structures:groups}{grup simetri} $\mathfrak{S}_3$
atas $\{1,2,3\}$ (dengan enam bijeksi: identitas, tiga transposisi, dua
siklus-$3$), lalu tunjukkan dua unsur yang tidak berkomutasi.
\end{exercise}

\begin{exercise}[$\star$]\label{exo:b1:structures:4}
Buktikan bahwa $H = \{z \in \C^* : \abs z = 1\}$ adalah \omterm{def:b1:structures:subgroup}{subgrup} $(\C^*,
\times)$, dan bahwa $\R_+^*$ juga demikian; apakah $H \cup \R_+^*$ sebuah
\omterm{def:b1:structures:subgroup}{subgrup}?
\end{exercise}

\begin{exercise}[$\star\star$]\label{exo:b1:structures:5}
Misalkan $f \colon (\R, +) \to (\C^*, \times)$, $\theta \mapsto
\eu^{\iu\theta}$. Buktikan bahwa $f$ \omterm{def:b1:structures:morphism}{morfisma}, hitung $\ker f$ dan
$\operatorname{im} f$, lalu simpulkan dari
\cref{prop:b1:structures:kernel} bahwa $f$ tidak \omterm{def:b1:logic:inj}{injektif}. Batasi
daerah asalnya agar ia \omterm{def:b1:logic:inj}{injektif} pada selang selebar mungkin.
\end{exercise}

\begin{exercise}[$\star\star$]\label{exo:b1:structures:6}
Misalkan $H, K$ dua \omterm{def:b1:structures:subgroup}{subgrup} $G$. Buktikan bahwa $H \cap K$ \omterm{def:b1:structures:subgroup}{subgrup},
dan bahwa $H \cup K$ menjadi \omterm{def:b1:structures:subgroup}{subgrup} \emph{hanya} bila $H \subseteq K$ atau
$K \subseteq H$. \emph{(Jika $h \in H \setminus K$ dan $k \in K
\setminus H$, di mana $hk$ dapat hidup?)}
\end{exercise}

\begin{exercise}[$\star\star$]\label{exo:b1:structures:7}
Sebuah \omterm{def:b1:structures:group}{grup} $G$ memenuhi $x^2 = e$ untuk setiap $x \in G$. Buktikan bahwa $G$
\omterm{def:b1:structures:group}{abelian}. \emph{(Jabarkan $(xy)^2$.)}
\end{exercise}

\begin{exercise}[$\star\star$]\label{exo:b1:structures:8}
Di $\Z/18\Z$: daftarkan unitnya lalu carilah invers $\overline 5$;
selesaikan $\overline 5\, x = \overline 7$; selesaikan $\overline 6\, x =
\overline 3$ dan $\overline 6\, x = \overline{12}$.
\end{exercise}

\begin{exercise}[$\star\star$]\label{exo:b1:structures:9}
Buktikan bahwa \omterm{def:b1:logic:sets}{himpunan} $\Z[\sqrt 2] = \{a + b\sqrt 2 : a, b \in \Z\}$ adalah
sebuah \omterm{def:b1:structures:ring}{ring} (subring $\R$), dan bahwa $1 + \sqrt 2$ merupakan unitnya
yang mempunyai tak hingga banyak pangkat berbeda --- sehingga $\Z[\sqrt 2]^\times$
tak hingga, tidak seperti $\Z^\times$.
\end{exercise}

\begin{exercise}[$\star\star\star$]\label{exo:b1:structures:10}
(\omterm{def:b1:structures:ring}{Ring} Boole) Misalkan $A$ sebuah \omterm{def:b1:structures:ring}{ring} yang di dalamnya $x^2 = x$ untuk setiap $x$.
Buktikan bahwa $x + x = 0$ untuk setiap $x$, dan bahwa $A$ komutatif.
\emph{(Jabarkan $(x+x)^2$ dan $(x+y)^2$.)} Berikan contoh \omterm{def:b1:structures:ring}{ring}
semacam itu dengan $\mathcal{P}(E)$, dengan mengambil selisih simetris sebagai
penjumlahan dan irisan sebagai perkalian.
\end{exercise}

\begin{exercise}[$\star\star\star$]\label{exo:b1:structures:11}
Misalkan $G$ sebuah \omterm{def:b1:structures:group}{grup} yang di dalamnya, untuk suatu $n \geq 1$ yang tetap, $(xy)^n =
x^n y^n$, $(xy)^{n+1} = x^{n+1}y^{n+1}$ dan $(xy)^{n+2} =
x^{n+2}y^{n+2}$ untuk setiap $x, y$. Buktikan bahwa $G$ \omterm{def:b1:structures:group}{abelian}.
\emph{(Dari ketiga kesamaan itu, turunkan lebih dulu $y^n x = x y^n$, lalu
$y^{n+1} x = x y^{n+1}$, dan simpulkan.)}
\end{exercise}

\begin{exercise}[$\star\star$]\label{exo:b1:structures:12}
\begin{enumerate}
  \item Tentukan semua \omterm{def:b1:structures:morphism}{morfisma} \omterm{def:b1:structures:group}{grup} dari $(\Z, +)$ ke $(\Z, +)$.
  \item Buktikan bahwa satu-satunya \omterm{def:b1:structures:morphism}{morfisma} \omterm{def:b1:structures:group}{grup} dari $(\Q, +)$ ke
        $(\Z, +)$ adalah \omterm{def:b1:structures:morphism}{morfisma} nol. \emph{(Untuk $x \in \Q$ dan
        $n \in \N^*$, bandingkan $f(x)$ dan $n\,f(x/n)$.)}
\end{enumerate}
\end{exercise}

\section{Soal: Grup simetri dan teka-teki 8}

\begin{problem}\label{pb:b1:structures:1}
\omterm{def:b1:structures:group}{Grup} $\mathfrak S_n$ berisi \omterm{def:b1:counting:objects}{permutasi} $\intint1n$ adalah
\omterm{def:b1:structures:group}{grup} tertua dalam matematika dan tetap yang paling mendidik. Soal
ini membangun teori strukturnya dari nol --- siklus,
pembangkitan oleh transposisi, \omterm{def:b1:structures:morphism}{morfisma} \emph{tanda}
$\varepsilon \colon \mathfrak S_n \to \{\pm1\}$ (yang keberadaannya
sungguh-sungguh tak sepele), dan \omterm{pb:b1:structures:1}{grup alternating} $\mathfrak A_n$
yang dibangkitkan oleh siklus-$3$ --- lalu menguangkannya pada sebuah
teka-teki klasik: pada permainan ubin geser $3 \times 3$, tak ada rangkaian langkah
yang dapat menukarkan dua ubin sambil membiarkan yang lain di tempatnya.
\index{grup simetri}\index{tanda}\index{grup alternating}
\omterm{def:b1:counting:objects}{Permutasi} bekerja pada $\intint1n$; hasil kali $\sigma\tau$ berarti
``terapkan $\tau$ lebih dulu''; dan $[\,v_1, \dots, v_n]$ menyatakan
\omterm{def:b1:counting:objects}{permutasi} yang mengirim $i$ ke $v_i$.

\medskip
\noindent\textbf{Bagian I --- Siklus dan transposisi.}
\begin{enumerate}
  \item Berikan alasan bahwa $\abs{\mathfrak S_n} = n!$
        (\cref{thm:b1:counting:counts}). Di $\mathfrak S_3$,
        hitung kedua hasil kali $\sigma = [2, 3, 1]$ dan $\tau =
        [1, 3, 2]$, lalu simpulkan bahwa $\mathfrak S_3$ tidak
        \omterm{def:b1:structures:group}{abelian}.
  \item Sebuah \emph{siklus-$k$} $(a_1\ a_2\ \dots\ a_k)$ ($k \geq 2$,
        dengan $a_i$ berbeda dua-dua) mengirim $a_1 \mapsto a_2
        \mapsto \dots \mapsto a_k \mapsto a_1$ dan menetapkan
        selebihnya; \emph{tumpuannya} adalah $\{a_1, \dots,
        a_k\}$. Buktikan bahwa dua siklus yang tumpuannya saling lepas
        berkomutasi.
  \item Buktikan bahwa setiap $\sigma \in \mathfrak S_n$ merupakan hasil kali
        siklus yang tumpuannya saling lepas dua-dua, dan bahwa
        penguraian itu tunggal sampai urutan faktornya.
        \emph{(Tinjau, untuk masing-masing $i$, barisan $i, \sigma(i),
        \sigma^2(i), \dots$: ia pasti kembali ke $i$; dan
        \emph{orbit} yang dihasilkannya mempartisi $\intint1n$, sedangkan $\sigma$ bekerja pada
        masing-masingnya sebagai siklus.)}
  \item Uraikan $\sigma = [4, 1, 5, 2, 3, 7, 8, 6] \in
        \mathfrak S_8$ menjadi siklus yang saling lepas. Dengan mendefinisikan
        \emph{\omterm{def:b1:structures:order}{orde}} $\sigma$ seperti pada
        \cref{def:b1:structures:order}, buktikan bahwa \omterm{def:b1:structures:order}{orde} sebuah
        hasil kali siklus yang saling lepas adalah KPK panjangnya,
        lalu hitung \omterm{def:b1:structures:order}{orde} $\sigma$ ini.
  \item Buktikan kesamaan teleskopis
        \[
        (a_1\ a_2\ \dots\ a_k)
        = (a_1\ a_k)(a_1\ a_{k-1})\cdots(a_1\ a_2) ,
        \]
        lalu simpulkan bahwa setiap \omterm{def:b1:counting:objects}{permutasi} merupakan hasil kali
        transposisi. Tuliskan $\sigma$ pada pertanyaan 4 sebagai hasil
        kali semacam itu.
  \item Tunjukkan lebih jauh bahwa transposisi \emph{bersebelahan}
        $(i\ \ i{+}1)$ sudah cukup: untuk $a < b$,
        \[
        (a\ b) = (a\ \ a{+}1)(a{+}1\ \ a{+}2)\cdots(b{-}1\ \ b)
        \cdots(a{+}1\ \ a{+}2)(a\ \ a{+}1),
        \]
        yaitu hasil kali $2(b - a) - 1$ transposisi bersebelahan ---
        yakni bilangan yang \emph{ganjil} (paritas ini akan penting dua kali
        di bawah).
\end{enumerate}

\medskip
\noindent\textbf{Bagian II --- Tandanya ada.} Untuk $\sigma
\in \mathfrak S_n$, misalkan
\[
N(\sigma) = \#\bigl\{(i, j) : i < j,\ \sigma(i) >
\sigma(j)\bigr\}
\]
banyaknya \emph{inversi} pada \omterm{def:b1:counting:objects}{permutasi} itu, lalu tulis $\varepsilon(\sigma) =
(-1)^{N(\sigma)}$.
\begin{enumerate}[resume]
  \item Hitung $N$ dan $\varepsilon$ untuk identitasnya, untuk sebuah
        transposisi $(i\ \ i{+}1)$, dan untuk $[2, 3, 1]$.
  \item Buktikan bahwa untuk setiap $\sigma$ dan setiap transposisi
        bersebelahan $\tau = (i\ \ i{+}1)$ berlaku
        $N(\sigma\tau) = N(\sigma) \pm 1$. \emph{(Mengomposisikan dengan
        $\tau$ di sebelah kanan menukarkan nilai pada posisi $i$ dan
        $i + 1$; dan tepat satu pasangan berubah status
        inversinya.)}
  \item Simpulkan, dengan memakai pertanyaan 6, bahwa untuk transposisi
        $\tau$ yang \emph{mana pun} berlaku $\varepsilon(\sigma\tau) =
        -\varepsilon(\sigma)$; lalu simpulkan bahwa jika $\sigma$
        hasil kali $p$ transposisi, maka $\varepsilon(\sigma)
        = (-1)^p$ --- khususnya paritas $p$ hanya bergantung
        pada $\sigma$, bukan pada pemfaktoran yang dipilih --- dan bahwa
        $\varepsilon \colon \mathfrak S_n \to \{\pm 1\}$ merupakan
        \omterm{def:b1:structures:morphism}{morfisma} \omterm{def:b1:structures:group}{grup}.
  \item Tunjukkan bahwa siklus-$k$ bertanda $(-1)^{k-1}$, dan
        bahwa secara umum $\varepsilon(\sigma) = (-1)^{n -
        c(\sigma)}$, dengan $c(\sigma)$ banyaknya orbit
        $\sigma$ (termasuk titik tetapnya).
  \item \emph{\omterm{pb:b1:structures:1}{Grup alternating}} adalah $\mathfrak A_n =
        \ker\varepsilon$. Berikan alasan bahwa ia \omterm{def:b1:structures:subgroup}{subgrup}, lalu buktikan
        $\abs{\mathfrak A_n} = \frac{n!}2$ untuk $n \geq 2$.
        \emph{(Tetapkan sebuah transposisi $\tau_0$ lalu tinjau $\sigma
        \mapsto \sigma\tau_0$.)}
  \item Pemeriksaan kesejalanan pada $\sigma = [4, 1, 5, 2, 3, 7, 8, 6]$:
        hitung $\varepsilon(\sigma)$ dengan tiga cara --- dengan mencacah
        inversinya, dari tipe siklusnya lewat pertanyaan 10, dan dari
        cacah transposisi Anda pada pertanyaan 5.
\end{enumerate}

\medskip
\noindent\textbf{Bagian III --- $\mathfrak A_n$ dibangkitkan oleh
siklus-$3$.}
\begin{enumerate}[resume]
  \item Misalkan $a, b, c, d$ berbeda dua-dua. Periksa kedua
        kesamaan
        \[
        (a\ b)(a\ c) = (a\ c\ b),
        \qquad
        (a\ b)(c\ d) = (a\ c\ b)(a\ c\ d) .
        \]
  \item Buktikan bahwa untuk $n \geq 3$, setiap unsur $\mathfrak
        A_n$ merupakan hasil kali siklus-$3$. \emph{(\omterm{def:b1:counting:objects}{Permutasi}
        genap adalah hasil kali sejumlah genap
        transposisi; serap dua-dua sekaligus.)}
  \item Tuliskan $(1\ 2)(3\ 4)$ dan siklus-$5$ $(1\ 2\ 3\ 4\ 5)$
        secara eksplisit sebagai hasil kali siklus-$3$.
  \item Buktikan rumus konjugasinya: untuk setiap $\sigma \in
        \mathfrak S_n$,
        \[
        \sigma\,(a_1\ \dots\ a_k)\,\sigma^{-1}
        = \bigl(\sigma(a_1)\ \dots\ \sigma(a_k)\bigr) .
        \]
\end{enumerate}

\medskip
\noindent\textbf{Bagian IV --- Teka-teki 8.} Ubin $1, \dots, 8$
digeser di dalam bingkai $3 \times 3$ dengan satu sel kosong; sebuah \emph{langkah}
menggeser ubin yang bersebelahan dengan sel kosong ke dalamnya. Nomori selnya
$1, \dots, 9$ (baris demi baris; pada kedudukan terpecahkan ubin $i$ berada di
sel $i$ dan sel kosongnya di sel $9$). Perlakukan sel kosong sebagai
ubin kesembilan, sehingga sebuah kedudukan menjadi \omterm{def:b1:counting:objects}{permutasi} $\sigma \in \mathfrak
S_9$ (dengan ubin $\sigma(i)$ duduk di sel $i$).
\begin{enumerate}[resume]
  \item Tunjukkan bahwa sebuah langkah mengganti $\sigma$ dengan $\sigma \circ \tau$
        dengan $\tau$ transposisi kedua sel yang
        terlibat; lalu simpulkan bahwa setiap langkah membalik
        $\varepsilon(\sigma)$.
  \item Misalkan $d(\sigma)$ jarak taksi (baris ditambah
        kolom) antara sel kosong yang sekarang dan sel
        asalnya $9$. Tunjukkan bahwa setiap langkah mengubah $d$ sebesar $\pm1$,
        sehingga setiap langkah juga membalik $(-1)^{d(\sigma)}$. Simpulkan bahwa
        \[
        I(\sigma) = \varepsilon(\sigma)\cdot(-1)^{d(\sigma)}
        \]
        bersifat \emph{invarian} terhadap setiap langkah.
  \item Buktikan kemustahilan klasik teka-teki itu: kedudukan
        yang menukarkan ubin $7$ dan $8$ sambil membiarkan segalanya
        (termasuk sel kosongnya) di tempatnya tak dapat dicapai
        dari kedudukan yang terpecahkan.
  \item Kita terima konversnya (buktinya berupa induksi yang mendidik tetapi
        panjang): setiap kedudukan dengan $I = +1$ dapat dicapai. Simpulkan
        bahwa tepat separuh dari $8!$ kedudukan
        dengan sel kosong di tempat asalnya dapat dipecahkan, yakni
        $\frac{8!}2 = 20\,160$.
  \item Simpulkan dari pertanyaan 20 bahwa susunan ubin yang dapat
        dicapai dengan sel kosong di tempat asalnya tepat membentuk
        \omterm{def:b1:structures:subgroup}{subgrup} $\mathfrak A_8 \leq \mathfrak S_8$.
  \item Penerapan invariannya: dapatkah kita mencapai (a) kedudukan
        yang ubin $1, 2, 3$-nya terpermutasi secara siklik
        sedangkan selebihnya, termasuk sel kosongnya, di tempat asalnya?
        (b) kedudukan yang ubin $5$ dan sel kosongnya sudah
        bertukar tempat sedangkan semua ubin lainnya di tempat asalnya? Berikan alasan
        bagi kedua jawabannya dengan $I$.
\end{enumerate}

\medskip
\noindent\textbf{Bagian V --- Sintesis.}
\begin{enumerate}[resume]
  \item Buktikan bahwa untuk $n \geq 3$ satu-satunya \omterm{def:b1:structures:morphism}{morfisma} \omterm{def:b1:structures:group}{grup} $f
        \colon \mathfrak S_n \to \{\pm 1\}$ adalah \omterm{def:b1:structures:morphism}{morfisma} yang
        konstan dan $\varepsilon$. \emph{(Dengan memakai pertanyaan 16 dan
        kekomutatifan $\{\pm1\}$, tunjukkan bahwa $f$ bernilai sama
        pada semua transposisi.)}
  \item Di mana persisnya soal ini memakai: (i) gagasan \omterm{def:b1:structures:morphism}{morfisma}
        beserta \cref{prop:b1:structures:kernel}; (ii) asas
        pencacahan pada \cref{ch:b1:counting}; (iii) persoalan
        keterdefinisian yang diselesaikan pertanyaan 8--9? Satu
        kalimat untuk masing-masing.
  \item Sintesis, dalam satu paragraf pendek: satu fungsi paritas,
        yang dibuktikan terdefinisi dengan baik sekali saja, sekaligus menata
        struktur dalam $\mathfrak S_n$ (yaitu \omterm{def:b1:structures:subgroup}{subgrup}
        $\mathfrak A_n$), memutuskan sebuah teka-teki fisik, dan ---
        lewat rumus $\det = \sum_\sigma
        \varepsilon(\sigma)\cdots$ --- akan mendefinisikan determinan
        pada \cref{ch:b1:det}. Berilah komentar atas pola yang berulang ini:
        invarian mengubah ``coba semua rangkaian langkah'' menjadi satu
        perhitungan.

\end{enumerate}
\end{problem}
